\section{The unit clause on the abelian locus from unramified torsion}
\label{sec:psi-unit-source}
This section translates the height-one argument in
\srcid{PsiUnitOnOfFixedOverUnramifiedExtension} and the subsequent
\srcid{PsiUnitOnAbelianLocusUniformisation}, at the frozen source commit.
It proves the unit clause of N3 from the actual uniformisation and its
torsion identification.  It does not construct that uniformisation.
This distinction matters: good reduction of an arbitrary abelian variety
would not identify its torsion with the particular degeneration datum
used here.

\subsection{The objects and the precise integral coefficients}
Let \(R\) be a normal Noetherian domain of characteristic zero,
\(K=\operatorname{Frac}(R)\), and \(F=\overline K\).
The degeneration datum \(D\) has a free character lattice \(X\), an
abelian scheme \(B/R\), an invertible sheaf \(L\) on \(B\), its
rigidification, symmetry and cubical generators, integral sections
\(c_\chi:\operatorname{Spec}R\to B\) for \(\chi\in X\), and generic
generators \(\psi_\chi\) of \(c_\chi^*L\).
All the structural generators in this sentence are the specified
generators of \(D\); replacing one by an arbitrary rational multiple
would change the question.

Write \(\widetilde G_D(F)\) for the group of points of the extension
attached to the cubical datum, and \(\iota:X\to\widetilde G_D(F)\) for
its lattice map.  A level-\(k\) point is a pair
\[
          p=(y,g),\qquad y\in X,\quad g\in\widetilde G_D(F),
          \qquad [k]g=\iota(y).
\]
Its equivalence is the actual lattice equivalence:
translation by \(a\in X\) changes the lattice label by \(ka\) and the
extension point by \(\iota(a)\), with the simultaneous signs determined
by the chosen direction of the equivalence.  In particular, an
equivalence between two points with the same label forces \(ka=0\).
The lattice \(X\) is torsion free.

Choose for every point \(b\in B\) an index \(J_0(b)\) of a chart of
\(L\) containing \(b\).  If \(P\) is a prime of \(R\), define
\[
 a_P(\chi)=\psiLoc(D,\chi,J_0(c_\chi(P)))\in K^\times.
 \tag{PU1}\label{eq:psi-local-coefficient}
\]
This is the generic coefficient of \(\psi_\chi\) in the indicated
integral chart.  The desired condition
\(\operatorname{PsiUnitOn}(D,\chi,U)\) says more than the existence of a
unit in one unrelated frame: in every integral trivialisation on every
open set contained in \(U\), the corresponding generic coefficient
extends to an invertible regular function.

\begin{srlemma}[A coefficient which is a local unit gives an actual neighbourhood]
\label{lem:psi-near-unit}
If \(a_P(\chi)=r/s\) with \(r,s\in R\setminus P\), then there is an
open neighbourhood \(V\) of \(P\) on which
\(\operatorname{PsiUnitOn}(D,\chi,V)\) holds.
\end{srlemma}
\begin{proof}
Intersect \(D(r)\), \(D(s)\), and
\(c_\chi^{-1}(U_{J_0(c_\chi(P))})\).  In this chart \(r/s\) is an
invertible regular function.  On an intersection with another chart,
the coefficient is obtained by multiplication by the corresponding
integral transition unit, or its inverse according to the coefficient
convention.  Thus it too is an invertible regular function.
The two expressions have the same generic value, so they are precisely
the coefficients of the original generic generator.
Restriction to any smaller open set preserves this equality.  This is
the assertion with all chart indices quantified, as in
\srcid{psiUnitOn_near}.
\end{proof}

\subsection{Integral evaluation of the cubical identities}
Fix a height-one prime \(P\).  Then \(A=R_P\) is a discrete valuation
ring; let \(\pi\) be a uniformiser.  For each \(\chi\), write uniquely
\[
             a_P(\chi)=u_\chi\pi^{m(\chi)},\qquad
             u_\chi\in A^\times,\quad m(\chi)\in\mathbb Z .
 \tag{PU2}\label{eq:psi-valuations}
\]
Let \(T\subset F\) be a discrete valuation ring containing \(A\) locally
such that the image of \(\pi\) is still a uniformiser.
Choose a valuation ring \(W\subset F\) dominating \(T\).  Its additive
valuation takes values in an ordered abelian group \(\Gamma\).
Set \(\alpha=v_W(\pi)>0\); then the map
\[
                   \mathbb Z\longrightarrow\Gamma,\qquad
                   j\longmapsto j\alpha
 \tag{PU3}\label{eq:psi-injective-value}
\]
is injective.

Here is the frame bookkeeping needed to evaluate the identities.
Every \(F\)-point of \(B\) extends to a \(W\)-point, by properness.
If two such extensions have the same generic point, separatedness
makes them equal.  We may therefore choose the chart of \(L\) at an
\(F\)-point to be \(J_0\) evaluated at the closed point of its integral
extension.  This gives a consistent chart choice for every integral
point simultaneously.  In any tensor combination of pullbacks of \(L\),
the coefficient of an integral generating section is a unit of \(W\).
Indeed, every chosen local frame is an integral basis, so the
coefficient of a generating section in their tensor product is
invertible.  Negative tensor powers use the inverse unit.  Evaluation
in \(F\) commutes with this integral evaluation.  These are exactly
\srcid{exists_good_ext} and \srcid{val_integral}; no boundedness
assertion about arbitrary rational frames is being used.

For \(g\in\widetilde G_D(F)\) and \(\chi\in X\), denote its coefficient
in this consistent chart system by \(z_\chi(g)\in F^\times\).
The cubical addition law has the form
\[
 z_\chi(g+h)
  =z_\chi(g)z_\chi(h)
     \operatorname{th}(c_\chi,g_B,h_B),
\]
where \(\operatorname{th}\) is the coefficient of the specified
cubical generator, and \(g_B,h_B\) are the points of \(B\).
All three points in this last evaluation are integral over \(W\).
Its valuation is therefore zero.  Induction, starting with the
rigidified value at zero, gives
\[
                  v_W(z_\chi([k]g))
                     =k\,v_W(z_\chi(g)).
 \tag{PU4}\label{eq:psi-nsmul-value}
\]
The specified biextension coefficient \(\tau(y,\chi)\) satisfies
\[
 \begin{split}
 v_W(\tau(y,\chi))
   &=v_W(a_P(y+\chi))-v_W(a_P(y))-v_W(a_P(\chi))\\
   &=\{m(y+\chi)-m(y)-m(\chi)\}\alpha .
 \end{split}
 \tag{PU5}\label{eq:psi-tau-value}
\]
To check the first equality, write the definition of the coefficient
as the product of the \(\psi_{y+\chi}\) coefficient, the inverses of the
\(\psi_y,\psi_\chi\) coefficients, and the inverse rigidification
coefficient.  The latter is an integral unit and has value zero.
Changing from the generic charts in this expression to the consistent
\(W\)-charts multiplies each coefficient by an integral transition
unit, also of value zero.  Thus the equality uses the very
coefficients in \eqref{eq:psi-valuations}.  This is the source's
\srcid{w_tau}, not a new pairing chosen to have that valuation.

Finally, the symmetry relation \(\operatorname{PsiEven}\) identifies
the coefficients at \(-\chi\) and \(\chi\) up to the coefficient of the
specified integral symmetry generator at \(c_\chi\).
Consequently
\[
                         m(-\chi)=m(\chi).
 \tag{PU6}\label{eq:psi-even-value}
\]
This follows from equality of their values in \(\Gamma\) and the
injectivity in \eqref{eq:psi-injective-value}.  The same calculation is
\srcid{w_even}.  The identities just proved also apply to the integral
extension of any multiple of \(g_B\), which is why
\eqref{eq:psi-nsmul-value} can be used for arbitrary \(k\).

\begin{srlemma}[The fixed-field step]
\label{lem:psi-fixed-field}
If \(x\in F^\times\) is fixed by every automorphism of \(F\) fixing
\(T\), then
\[
                          x=u\pi^v
                  \quad (u\in T^\times,\ v\in\mathbb Z).
\]
\end{srlemma}
\begin{proof}
Let \(E=\operatorname{Frac}(T)\subset F\).
The field \(F\) is algebraic over \(E\), algebraically closed, and has
characteristic zero.  The fixed field of
\(\operatorname{Aut}(F/E)\) is therefore \(E\).
For completeness, if an algebraic \(x\notin E\) were fixed, its
separable minimal polynomial would have another root in \(F\).
The \(E\)-embedding of \(E(x)\) sending \(x\) to that root extends to an
automorphism of the algebraic closure, a contradiction.
Thus \(x\in E^\times\), where the discrete valuation decomposition is
exactly a unit times the integral power of the chosen uniformiser.
\end{proof}

\begin{srlemma}[One level gives divisibility]
\label{lem:psi-level-divisibility}
Suppose that \(k\geq1\), and every automorphism of \(F\) fixing \(T\)
fixes every level-\(k\) point up to the actual lattice equivalence.
Then
\[
       k\mid m(y+\chi)-m(y)-m(\chi)\qquad(y,\chi\in X).
 \tag{PU7}\label{eq:psi-level-divisibility}
\]
\end{srlemma}
\begin{proof}
Choose \(p=(y,g)\) with \([k]g=\iota(y)\).
Such a point exists in the source by the multiplication law
\(\operatorname{MulNSt}\).  On the actual extension its elementary
content is as follows: multiplication by \(k\) is surjective on the
abelian quotient over the algebraically closed field \(F\).
After lifting a preimage there, the discrepancy is a point of the
torus; multiplication by \(k\) is surjective on the split torus since
every nonzero element of \(F\) has a \(k\)-th root.  Correcting by that
torus point gives the required \(g\).

For an automorphism \(\sigma\) fixing \(T\), the labels of \(p\) and
\(\sigma p\) are both \(y\).  The asserted lattice equivalence is
therefore represented by \(a\in X\) with \(ka=0\).
Since \(X\) is torsion free, \(a=0\), and hence
\(\sigma g=g\) as an actual extension point.
The coefficient \(z_\chi(g)\) is now fixed by \(\sigma\).
Indeed, its defining four pullbacks of \(L\) are at
\(c_\chi+g_B,c_\chi,g_B,0\); these maps and their chosen chart indices
are fixed.  Functoriality of evaluation under the field automorphism
then fixes the coefficient itself.  This verifies
\srcid{charVal_fixed}, including the frame issue.

Lemma~\ref{lem:psi-fixed-field} gives
\(z_\chi(g)=u\pi^v\), with \(u\in T^\times\).
The lattice embedding has coefficient
\(z_\chi(\iota(y))=\tau(y,\chi)\), in the specified cubical convention.
Apply valuations to \([k]g=\iota(y)\), and use
\eqref{eq:psi-nsmul-value} and \eqref{eq:psi-tau-value}.
The result is
\[
       kv\alpha
          =\{m(y+\chi)-m(y)-m(\chi)\}\alpha .
\]
Injectivity of \eqref{eq:psi-injective-value} gives the required
integer divisibility.  Notice why the extension is unramified in
exactly the required sense: the uniformiser remains \(\pi\).
An arbitrary ramified extension would introduce a varying
ramification factor and would not prove \eqref{eq:psi-level-divisibility}.
\end{proof}

\begin{srtheorem}[The height-one unit statement]
\label{thm:psi-height-one-unit}
Suppose there is \(n\geq2\) such that, for every \(j\geq0\), a discrete
valuation subring \(T_j\subset F\) as above leaves \(\pi\) a uniformiser
and makes every level-\(n^j\) point fixed up to lattice equivalence.
Then \(a_P(\chi)\in R_P^\times\) for every \(\chi\), and each
\(\psi_\chi\) is a unit on a neighbourhood of \(P\).
\end{srtheorem}
\begin{proof}
Fix \(y,\chi\), and put
\(\delta=m(y+\chi)-m(y)-m(\chi)\).
Lemma~\ref{lem:psi-level-divisibility} shows \(n^j\mid\delta\) for
every \(j\).  Choose \(j\) with \(n^j>|\delta|\); a nonzero multiple of
\(n^j\) has absolute value at least \(n^j\), so \(\delta=0\).
Thus \(m\) is additive.  Additivity at \(0+0\) gives \(m(0)=0\).
Additivity at \(\chi+(-\chi)\), combined with
\eqref{eq:psi-even-value}, gives \(2m(\chi)=0\), whence \(m(\chi)=0\)
in \(\mathbb Z\).
Equation~\eqref{eq:psi-valuations} proves the first assertion.
Writing this local unit as \(r/s\), with both \(r,s\notin P\), and
using Lemma~\ref{lem:psi-near-unit} proves the second.
This is the argument of \srcid{psiUnitOn_of_unramified}; no uniform
choice of \(T_j\) in \(j\) is needed.
\end{proof}

\subsection{The source's unramified torsion producer}
The producer \srcid{exists_unramified_torsion_points} uses lifting of
finitely many special-fibre points and equality of torsion counts.
Here are those steps, rather than an assumption that all torsion is
already unramified.

\begin{srlemma}[Simultaneous integral lifts of an etale finite set]
\label{lem:psi-etale-finite-lifts}
Let \(A\) be a discrete valuation ring, \(F\) an algebraically closed
field containing \(\operatorname{Frac}(A)\), and
\(\rho:A\to\Omega\) a local map to a field.
Let \(Z/A\) be etale and separated.  Given finitely many
\(\Omega\)-points of its special fibre, there are a discrete valuation
ring \(A'\), a local extension \(A\to A'\), an embedding \(A'\to F\),
and a compatible local map \(A'\to\Omega\), over which all those
points lift.  Every uniformiser of \(A\) is still a uniformiser of
\(A'\).
\end{srlemma}
\begin{proof}
First consider one point.  Choose compatible affine charts around it
and its image, so its coordinate map is a homomorphism
\(\epsilon:C\to\Omega\) from an etale \(A\)-algebra.
Localise \(C\) at \(\mathfrak q=\ker\epsilon\), obtaining \(A'\).
The map \(A'\to\Omega\) is local: a fraction whose image is nonzero
has numerator outside \(\mathfrak q\), and is a unit.
The map \(A\to A'\) is flat, local, essentially of finite type and
formally unramified.  Hence \(A'\) is a regular local ring of dimension
one and its maximal ideal is generated by the image of a uniformiser
of \(A\).  More explicitly, the zero-dimensional etale closed fibre
has a field as its local ring, so the maximal ideal is
\(\mathfrak m_AA'\); flatness preserves the nonzero-divisor
uniformiser, and the regular-local dimension formula gives dimension
one.  A one-dimensional regular local ring is a discrete valuation
domain.  Its fraction field is algebraic over
\(\operatorname{Frac}(A)\), by relative dimension zero.
It therefore embeds in \(F\) over that field; the embedding restricts
to the required embedding of \(A'\).
The coordinate homomorphism supplies the lift of the given point.

For a finite list, repeat this construction for the next point after
each local extension.  Pull back every previously obtained lift along
the new extension.  The new residue map remains compatible, the
generic-field embedding continues the old one, and uniformisers
remain uniformisers at every step.  Induction on the finite list
proves the assertion.  This is
\srcid{exists_lift_one} followed by \srcid{exists_lift_finset}.
\end{proof}

\begin{srlemma}[Equal point counts turn those lifts into all generic points]
\label{lem:psi-etale-count-lifts}
In Lemma~\ref{lem:psi-etale-finite-lifts}, assume the geometric special
and generic point sets are finite and have the same cardinality.
There is a discrete valuation subring \(T\subset F\), with the same
uniformiser as \(A\), over which every generic \(F\)-point of \(Z\)
extends.
\end{srlemma}
\begin{proof}
Lift the whole finite special point set simultaneously.
Send each special point to the generic point of its lift.
This map is injective: if two such generic points agree, their
sections over the integral ring \(A'\) agree on the dense generic
point, and separatedness makes the sections equal.  Their special
values are then equal.  An injection of finite sets of equal
cardinality is a bijection, so every generic point is one of these
lifted points.  Replace \(A'\) by its image under its embedding in
\(F\).  It is isomorphic to \(A'\), is a discrete valuation subring,
and has all the asserted properties.  This is
\srcid{exists_subalgebra_of_card}.
\end{proof}

\begin{srproposition}[Good reduction gives the required actual torsion lifts]
\label{prop:psi-good-reduction-torsion}
Let \(H/A\) be an abelian scheme over a discrete valuation ring, and
let \(k\) be a positive integer invertible in \(A\).
There is a discrete valuation subring \(T\subset F\), with the same
uniformiser as \(A\), over which every point of \(H_F[k]\) extends.
Every automorphism of \(F\) fixing \(T\) fixes these torsion points.
\end{srproposition}
\begin{proof}
The kernel \(Z=H[k]\) is etale and separated over \(A\).
For the etaleness, the differential of \([k]\) on the relative tangent
bundle is multiplication by the unit \(k\); translation transports
this calculation from the identity to every point.  The usual
smooth etale criterion gives etaleness, and its pullback to the
identity gives the assertion for the kernel.  Separatedness is
inherited from \(H\).
The relative dimension \(g\) is constant on this connected base.
The classical prime-to-characteristic torsion theorem for abelian
varieties gives exactly \(k^{2g}\) geometric torsion points in both
the generic and the special fibres.  This precise background theorem
is the source's
\srcid{card_torsion_points_abelianVariety}; no assertion about
\(p\)-power torsion in characteristic \(p\) is used.
Lemma~\ref{lem:psi-etale-count-lifts} applies.
A point which is the base change of a \(T\)-point is fixed by every
automorphism fixing \(T\), since its coordinate maps already have
coefficients in \(T\).  If the point is described through a Cartesian
generic-fibre diagram, equality can be checked after both
projections of that diagram, so the same conclusion holds for the
actual generic-fibre point.  This is the argument
\srcid{fixed_of_extends}.
\end{proof}

\subsection{Transport through the actual torsion identification}
For this step retain the actual uniformisation \(U\) and its actual
identification
\[
       Q_k(D,F)/{\sim}\ \longrightarrow\ H_F[k],
 \tag{PU8}\label{eq:psi-actual-torsion-id}
\]
whose equality fibres are precisely \(\sim\).
The identification must be the one defined by the theta action and
the expansion relation of \(U\).  A mere equality of cardinalities
between the two sets would not suffice.

The required naturality is the exact statement
\srcid{qkToAkU_iff_of_qkExt}.  Its scalar descent step can be stated
explicitly.  If \(V_E\to V_F\) is a genuine field extension of vector
spaces, operators \(S,T\) and \(S',T'\) commute with extension, and
\[
                         S'=cT'\quad(c\in F^\times),
\]
choose an \(E\)-linear functional \(\ell:F\to E\) with \(\ell(1)=1\)
and apply the induced coefficient retraction.  It gives
\(S=\ell(c)T\).  If \(\ell(c)\ne0\), this is the required
proportionality over \(E\).  If \(\ell(c)=0\), it gives \(S=0\);
injectivity of extension and \(c\ne0\) then give \(T=0\), so the
nonzero scalar \(1\) works.  Conversely proportionality over \(E\)
extends, because the extended images span \(V_F\).
This proves the source's \srcid{proportional_iff}, including its
otherwise easy-to-miss zero-operator case.

For the actual identification, the source applies this calculation
to the theta translation operator on algebraic sections and the
operator on expanded families.  Pullback of a line-bundle
isomorphism commutes with its action on sections; evaluation of the
pulled-back cubical coefficients is the pullback of their evaluation.
Consequently both operators commute with field extension.
The expansion equivalence commutes with extension as well, by
pullback of the actual section identity.  Its family extension is a
genuine extension of vector spaces: the original images span, and
the coefficient retractions are defined.  Proportionality is exactly
the theta-action definition of the relation in
\eqref{eq:psi-actual-torsion-id}, independent of the choice of a
generating frame since two such frames differ by a nonzero scalar.
These facts give, for an embedding \(\sigma:E\to F\), a level point
\(p'\) obtained from \(p\), and the corresponding actual points,
\[
 QkToAk_U(p,x)\ \Longleftrightarrow\
 QkToAk_U(p',\sigma x).
 \tag{PU9}\label{eq:psi-equivariance}
\]
This paragraph uses the actual expansion equivalence and its
base-change theorem; it does not construct that equivalence.
They remain separate uniformisation inputs until their own source
proofs have been translated.

\begin{srtheorem}[Unit clause on the abelian locus]
\label{thm:psi-unit-abelian-locus}
Let \(R\) now be complete local, normal Noetherian of characteristic
zero, with the source's admitted maximal-ideal data.
Let \(D\) be its actual degeneration datum and \(U\) an actual
uniformisation with the torsion identification and expansion
properties just specified.  Suppose the semiabelian family \(H/R\)
has open abelian locus \(V\).  Then, for every \(\chi\in X\),
\[
                         \operatorname{PsiUnitOn}(D,\chi,V).
\]
\end{srtheorem}
\begin{proof}
There is an integer \(n\in\{2,3\}\) which is a unit of \(R\).
Otherwise both \(3\) and \(-2\) would belong to the maximal ideal,
whereas their sum is \(1\).
Fix a height-one point \(P\in V\), and set \(A=R_P\).
Every point of \(\operatorname{Spec}A\) maps to a generalisation of
\(P\); since \(V\) is open, the entire localised family is an
abelian scheme.  For each \(j\), \(n^j\) is invertible in \(A\).
Proposition~\ref{prop:psi-good-reduction-torsion} supplies a
uniformiser-preserving \(T_j\subset F\) fixing every actual
\(n^j\)-torsion point.

Apply \eqref{eq:psi-equivariance} to an automorphism fixing \(T_j\).
The actual torsion image of a level point and that of its conjugate
are equal.  The equality-fibre assertion in
\eqref{eq:psi-actual-torsion-id} therefore makes the two level points
equivalent.  Theorem~\ref{thm:psi-height-one-unit} proves that
\(\psi_\chi\) is a unit near \(P\).

It remains to pass from height one to all of \(V\); this last step
uses normality and must be applied to both a coefficient and its
inverse.  Here is the local argument.  For \(a=b/c\in K\), \(c\ne0\),
let
\[
                    I=(c):(b)=\{s\in R:sb\in cR\}.
\]
The element \(a\) is regular at \(Q\) exactly when \(I\not\subset Q\).
If it is not regular, choose a prime \(\mathfrak p\subset Q\) minimal
over \(I\).  It is an associated prime of \(R/I\).
The injection
\[
                R/I\ \longrightarrow\ R/(c),\qquad
                \bar x\longmapsto\overline{xb},
\]
then makes \(\mathfrak p\) an associated prime of \(R/(c)\).
Every such prime has height one in a normal Noetherian domain:
if its height were at least two, Serre's \(S_2\) condition at
\(\mathfrak p\), followed by the depth formula for the nonzero divisor
\(c\), would give positive depth for \(R_{\mathfrak p}/(c)\),
contrary to being associated.  Its height is not zero since it
contains \(c\ne0\).
As \(V\) is open and \(Q\in V\), also \(\mathfrak p\in V\).
Regularity at this height-one prime supplies
\(s\notin\mathfrak p\) with \(sa\in R\), or \(s\in I\), contradicting
\(I\subset\mathfrak p\).
Thus any rational function regular at every height-one point of
\(V\) is regular at every point of \(V\).

Apply this argument on the integral trivialising opens for
\(c_\chi^*L\) to its coefficient \(a\), and again to \(a^{-1}\).
At every point both admit fractions with denominators outside that
point; after shrinking, these fractions are regular inverses.
They agree on overlaps because they agree in \(K\).
They therefore glue to an invertible regular coefficient on every
such open.  Lemma~\ref{lem:psi-near-unit}'s transition calculation
then gives the full chart-independent condition
\(\operatorname{PsiUnitOn}\).
This is the source's \srcid{psiUnitOn_of_height_one}, followed by
\srcid{psiUnitOn_abLocus}.
\end{proof}

\paragraph{Exact remaining boundary of this source translation.}
The proofs above discharge the height-one divisibility, local-unit,
unramified finite-point lifting, fixed-point transport and normal
extension arguments.  The standard background used in the finite
point count is the prime-to-characteristic torsion theorem for an
abelian variety, with its dimension \(g\) and invertibility condition
specified above; the normality background is Serre's criterion and
the nonzero-divisor depth formula.
The actual degeneration datum, its cubical laws, its expansion
equivalence and torsion equality-fibre theorem, and the construction
of \(U\) over a complete normal local base are not replaced by those
background facts.  Their construction remains part of the N3
translation frontier.  Consequently this section alone does not
supply all of N3 or the full integral-model package.

\paragraph{Source correspondence.}
All source paths below are at commit
\srcid{7cbf7a0c32fca1badf05931879a37c3f8a2b03b4}.
The valuation identities, fixed-field lemma, lattice-label
divisibility and integral neighbourhood construction are in
\srcid{Development/AbelianVarietyTheory/Proofs/DegenerationProofs/PsiUnitOnOfFixedOverUnramifiedExtension.lean}.
The finite special-point lifting, cardinality injection and
good-reduction producer are in
\srcid{Development/AbelianVarietyTheory/Proofs/DegenerationProofs/ExistsUnramifiedExtTorsionPointsGoodReduction.lean}.
The actual operator comparison is in
\srcid{Development/AbelianVarietyTheory/Proofs/UniformizationProofs/QkToAkUInvariantUnderFieldEmbedding.lean}.
The Hartogs step and the assembly over the abelian locus are in
\srcid{Development/AbelianVarietyTheory/Proofs/DegenerationProofs/PsiUnitOnAbelianLocusUniformisation.lean}.
